\chapter{Discussion}
\label{ch:discussion}

The preceding chapters reported four model comparisons in isolation, one per perception sub-task and one for the integrated system. Read together, the four comparisons share a few patterns worth stating directly, and they fail in a few places worth stating just as directly.

\section{Cross-Cutting Findings}
\label{sec:disc-findings}

The clearest pattern across the study is that, on both perception problems where a learned model could be set against a more direct geometric or rule-based alternative, the learned model won, and it won for the same underlying reason. For ball localisation (Section~\ref{sec:ball-analysis}), the heatmap formulation of TrackNetV4 was a far stronger localiser than the single-frame bounding-box regression of YOLO11m, which was the weakest model by a wide margin. For bounce detection (Section~\ref{sec:bounce-analysis}), both the gradient-boosted classifier and the temporal convolutional network separated genuine bounces from racket hits almost perfectly, even though a hit and a bounce produce the same vertical trajectory reversal and a naive kinematic rule keyed on that reversal cannot tell them apart. The common thread is that the discriminative information in each case lives in fine spatial or temporal shape rather than in a single cue. A tiny tennis ball spanning a handful of pixels is poorly served by a coarse bounding box, but well served by a dense Gaussian heatmap that the network learns to peak precisely \cite{huang2019tracknet, cheng2023smallobjectsurvey}; a bounce is poorly served by detecting a reversal, but well separated from a hit by the full local shape of the trajectory leading into and out of the contact. In both cases the learned representation captured shape that the simpler formulation discarded.

The second recurring finding is that model capacity was not the deciding factor. What mattered was whether the architecture carried the right inductive bias and produced output at the right resolution. The ball comparison is the sharper example: the V1--V5 progression was not monotone in accuracy, and TrackNetV4, whose lightweight motion-attention module injects an explicit motion prior \cite{raj2025tracknetv4}, outperformed the heavier TrackNetV5, whose learned temporal self-attention was left to discover the same cue from data. A strong, cheap inductive bias beat a general, expensive mechanism asked to learn that bias on its own. The court comparison (Section~\ref{sec:court-analysis}) makes the resolution half of the point: the smallest model considered, a MobileNetV3-Small backbone of 1.28 million parameters \cite{howard2019mobilenetv3}, was simultaneously the most accurate at strict tolerance and the fastest, while the two larger stride-four backbones were quantisation-limited and only became competitive once the tolerance was relaxed. Strict-tolerance accuracy there was governed by the heatmap output stride, not by the parameter count of the encoder. Taken together, the ball and court results support one claim: on these problems architectural fit to the task, the correct prior beats raw scale.

\section{Threats to Validity and Limitations}
\label{sec:disc-limitations}

Several factors bound how far these conclusions should be carried. The most concerning is test-set size, which is most acute for bounce detection. The held-out test was a single match (game7) of 1881 frames containing only 50 bounce events. At that scale each individual event is worth roughly two points of recall, and a handful of false positives moves event F1 by several points.

A second limitation concerns latency. The measured per-frame CPU latency of the composed pipeline does not support a real-time claim, and the integrated system is positioned accordingly as an offline analysis tool rather than a live one.

A third and qualitative limitation concerns the in/out line call produced by the integrated system (Section~\ref{sec:int-inout}). That call projects each visible bounce through the court homography into court-metre coordinates and tests polygon containment, but it is reported without any accuracy, precision, or error figure. No dataset carries ground-truth in/out labels against which it could be scored. The output should therefore be read as a plausible visual indication rather than an adjudication.

A fourth limitation is internal to the integrated pipeline. The bounce detector was scored in Chapter~\ref{ch:bounce} on the ground-truth ball trajectory, but in the assembled system (Section~\ref{sec:int-arch}) it consumes the TrackNetV4 trajectory instead. That predicted trajectory carries localisation error which, as the per-visibility breakdown showed, is largest precisely on the occluded near-ground frames where bounces occur. Deployed bounce timing therefore sits below the Chapter~\ref{ch:bounce} numbers by an amount this study did not isolate, and the reported bounce metrics should be read as an upper bound on end-to-end performance rather than a deployed figure.

Finally, the study deliberately omitted several capabilities that a complete match-analysis system would require: there is no player tracking, no pose estimation, no shot classification, no multi-camera fusion, and no 3D reconstruction. These omissions are scope choices rather than oversights, but they mean the system addresses only the ball, court, and bounce perception layer. Collectively, the limitations above motivate the future work outlined in Chapter~\ref{ch:conclusion}.
